\chapter{Décomposer la fourchette}\label{ch:mx:decomposing-the-spread}

Un achat de 1\,000 actions coûte deux centimes au-dessus du prix milieu. Une minute plus tard, le prix milieu a bougé d'un centime et demi dans le sens de la transaction~: le fournisseur de liquidité
a conservé un demi-centime, et le reste correspondait à ce que la transaction avait révélé au marché. Le chapitre~4 a expliqué pourquoi une fourchette existe~; le présent chapitre mesure ce que rémunère chacune
de ses parties. La mesure est plus difficile qu'il n'y paraît~: le mouvement qui suit une transaction mêle ce que la transaction a révélé, les transactions qui l'ont suivie et
le bruit propre du marché, et chaque méthode les démêle selon ses propres hypothèses. Nous testons chaque méthode sur \texttt{firm.tape}, le marché simulé
du livre~7, où le prix efficient est connu et où la sélection adverse réelle de chaque transaction peut être calculée.

\section{Fourchette affichée, effective et réalisée}

\begin{definition}[Fourchette affichée]\label{def:mx:decomposing-the-spread:quoted}
La \emph{fourchette affichée}\index{fourchette affichée} à l'instant $t$ vaut $a_t-b_t$, soit le meilleur cours vendeur moins le meilleur cours acheteur~; la fourchette affichée relative la divise par le
prix milieu $m_t=\tfrac12(a_t+b_t)$. Moyennée sur une période, elle est pondérée par le temps pendant lequel chaque cotation est restée en vigueur.
\end{definition}

La \omterm{def:mx:decomposing-the-spread:quoted}{fourchette affichée} est ce que paierait un petit ordre~; la fourchette effective (livre~1, chapitre~10) est ce que les transactions ont effectivement payé. Avec $\varepsilon_t=+1$
pour une transaction déclenchée par un acheteur et $-1$ pour une transaction déclenchée par un vendeur (le signe de la transaction, livre~7, chapitre~9), une transaction au prix $p_t$ a pour demi-fourchette effective
égale à $\varepsilon_t(p_t-m_t)$, mesurée par rapport au prix milieu qui la précède immédiatement. Elle est inférieure à la demi-fourchette affichée lorsque la transaction a reçu
une amélioration de prix, et supérieure lorsque l'ordre a parcouru le carnet. La fourchette réalisée mesure ce que le fournisseur de liquidité a conservé une fois le prix déplacé.

\begin{definition}[Impact de la transaction sur le prix]\label{def:mx:decomposing-the-spread:impact}
L'\emph{impact de la transaction sur le prix}\index{impact de la transaction sur le prix} d'une transaction à l'horizon $h$ vaut $\varepsilon_t(m_{t+h}-m_t)$~: le mouvement du prix milieu dans le sens de la
transaction sur les $h$ suivants. La demi-fourchette réalisée vaut $\varepsilon_t(p_t-m_{t+h})$, ce que le fournisseur gagnerait en clôturant sa position au
prix milieu ultérieur, et transaction par transaction~:
\[
\underbrace{\varepsilon_t(p_t-m_t)}_{\text{effective}}=\underbrace{\varepsilon_t(p_t-m_{t+h})}_{\text{réalisée}}+\underbrace{\varepsilon_t(m_{t+h}-m_t)}_{\text{impact}}.
\]
\end{definition}

\begin{omfigure}
\begin{tikzpicture}[font=\footnotesize,x=1cm,y=1.1cm]
\draw[->,gray] (0,0) -- (9.5,0) node[right] {temps};
\draw[omDef,very thick] (0.3,1) -- (3,1) node[pos=0.1,above] {prix milieu $m_t$};
\draw[omDef,very thick] (3,1) -- (4.2,1.5) -- (5.4,1.9) -- (7,2.5) -- (9.2,2.5);
\node[omDef,right] at (9.2,2.5) {prix milieu $m_{t+h}$};
\fill[omThm] (3,3) circle (2.5pt) node[above] {achat à $p_t$};
\draw[dashed,gray] (3,0) -- (3,3); \node[below] at (3,0) {$t$};
\draw[dashed,gray] (7.8,0) -- (7.8,2.5); \node[below] at (7.8,0) {$t+h$};
\draw[dotted] (3,3) -- (9.2,3);
\draw[decorate,decoration={brace,amplitude=4pt,mirror}] (2.8,1) -- (2.8,3) node[midway,left=5pt,align=right] {demi-fourchette\\ effective};
\draw[decorate,decoration={brace,amplitude=4pt}] (8.1,3) -- (8.1,2.5) node[midway,right=5pt] {réalisée};
\draw[decorate,decoration={brace,amplitude=4pt}] (8.1,2.5) -- (8.1,1) node[midway,right=5pt] {impact};
\draw[dotted] (3,1) -- (8.1,1);
\end{tikzpicture}

\omcaption{Un achat, scindé en deux. La demi-fourchette effective est ce que l'acheteur a payé au-dessus du prix milieu~; après $h$, la part dont le prix milieu s'est déplacée dans le sens de la
transaction constitue l'impact, et la part qui revient au fournisseur de liquidité est la fourchette réalisée.}\label{fig:mx:decomposing-the-spread:split}
\end{omfigure}

\omcode{code/firm/spreaddecomp/firm_spreaddecomp.py}{41}{48}{Demi-fourchettes effectives, réalisées et d'impact pondérées par les actions, à un horizon donné.}\label{lst:mx:decomposing-the-spread:decompose}

La session simulée dure 6.5 heures, compte 21\,993 transactions de 118 actions en moyenne, et sa \omterm{def:mx:decomposing-the-spread:quoted}{fourchette affichée} pondérée par le temps vaut 1.04 tick. La demi-fourchette effective des transactions
pondérée par le nombre d'actions comme le font les rapports réglementaires, vaut 0.535 tick.

\begin{dated}{2026-09}{Fourchettes réalisées dans les rapports Rule 605}\label{dat:mx:decomposing-the-spread:rule605}
Les amendements de la SEC à la Rule 605 de la Regulation NMS (Release 34-99679, adoptés le 6 mars 2024, publiés le 15 avril 2024) exigent que les rapports sur la qualité d'exécution
donnent les fourchettes réalisées à cinq horizons, 50 millisecondes, 1 seconde, 15 secondes, 1 minute et 5 minutes, au lieu de l'unique
horizon de 5 minutes de la règle initiale, et que les horaires soient mesurés en millisecondes ou plus finement. La date de mise en conformité, fixée au 14 décembre 2025, a été reportée le
30 septembre 2025 au premier août 2026.
\end{dated}

\section{L'impact de la transaction sur le prix}

\begin{definition}[Composante de sélection adverse]\label{def:mx:decomposing-the-spread:adverse}
La \emph{composante de sélection adverse}\index{composante de sélection adverse} de la fourchette effective est la part qui rémunère les fournisseurs de liquidité pour
avoir traité avec des contreparties qui en savent davantage~: pour une transaction à l'instant $t$, il s'agit de l'espérance de $\varepsilon_t(P^\ast_t-m_t)$, où $P^\ast_t$ est le prix efficient,
c'est-à-dire l'espérance de la valeur de l'actif compte tenu de toute l'information, y compris ce que sait l'agresseur et que les cotations ne montrent pas encore.
\end{definition}

\begin{definition}[Décomposition de la fourchette]\label{def:mx:decomposing-the-spread:decomposition}
Une \emph{décomposition de la fourchette}\index{décomposition de la fourchette} répartit la fourchette effective entre une \omterm{def:mx:decomposing-the-spread:adverse}{composante de sélection adverse} et le reste, c'est-à-dire
les \omterm{def:mx:why-there-is-a-spread:processing}{coûts de traitement des ordres} et les \omterm{def:mx:why-there-is-a-spread:inventory}{coûts de détention de l'inventaire} décrits au chapitre~4, que le fournisseur conserve une fois le prix stabilisé. Les méthodes diffèrent par ce qu'elles observent
(transactions seules, transactions et cotations, barres journalières) et par le modèle qui leur permet de distinguer les deux parts.
\end{definition}

Sur un marché réel, $P^\ast$ n'est pas observé, et la méthode de la fourchette réalisée le remplace par un prix milieu ultérieur. Elle fonctionne lorsque le prix milieu rejoint le
prix efficient et que rien d'autre ne le déplace dans le sens de la transaction.

\begin{proposition}[Vers quoi converge l'impact]\label{prop:mx:decomposing-the-spread:converge}
Soit $P^\ast$ une martingale, soient $\varepsilon_t$ et $m_t$ connus à l'instant $t$, et supposons que $\E[\varepsilon_t(m_{t+h}-P^\ast_{t+h})]$ tende vers zéro quand $h$ croît. Alors
\[
\E[\varepsilon_t(m_{t+h}-m_t)]\;\longrightarrow\;\E[\varepsilon_t(P^\ast_t-m_t)],
\]
c'est-à-dire la \omterm{def:mx:decomposing-the-spread:adverse}{composante de sélection adverse}. Le mouvement propre du prix efficient ajoute à l'impact de chaque transaction un terme d'espérance nulle et de variance $\sigma^2h$, où
$\sigma^2$ est le taux de variance de $P^\ast$~: sur $N$ transactions indépendantes, sa moyenne a pour écart type $\sigma\sqrt{h/N}$.
\end{proposition}

\begin{proof}
Écrivons $\varepsilon_t(m_{t+h}-m_t)=\varepsilon_t(m_{t+h}-P^\ast_{t+h})+\varepsilon_t(P^\ast_{t+h}-P^\ast_t)+\varepsilon_t(P^\ast_t-m_t)$. Le terme central est
d'espérance nulle~: $\varepsilon_t$ est connu à l'instant $t$ et les accroissements de $P^\ast$ après $t$ sont d'espérance conditionnelle nulle. Le premier terme tend vers zéro par hypothèse. Le terme central
a pour variance $\E[\varepsilon_t^2(P^\ast_{t+h}-P^\ast_t)^2]=\sigma^2h$, puisque $\varepsilon_t^2=1$, et la moyenne de $N$ copies indépendantes a pour variance
$\sigma^2h/N$.
\end{proof}

La proposition traduit un compromis. Un horizon court n'a pas laissé au prix milieu le temps de rattraper le prix, et il prend pour fourchette gagnée ce qui relève de l'information~; un horizon long laisse le
mouvement propre du prix efficient noyer la moyenne. Sur un marché réel, les transactions se chevauchent aussi~: l'impact d'une transaction à 5 minutes contient l'impact de toutes les transactions survenues pendant ces
5 minutes, y compris le reste de son propre métaordre.

La \Cref{fig:mx:decomposing-the-spread:horizon} montre ces deux effets sur la journée simulée. La part de sélection adverse réelle vaut 0.449 tick, soit 84~\% de la
demi-fourchette effective~: les \omterm{def:mx:why-there-is-a-spread:informed}{traders informés}, qui représentent 21~\% du volume, négocient lorsque le prix efficient se situe à 2.05 ticks au-delà du prix milieu de leur côté, tandis que les \omterm{def:mx:why-there-is-a-spread:informed}{noise traders} le font lorsque
le prix efficient se situe à 0.025 tick au-delà. L'impact monte de 0.03 tick à 50 millisecondes à 0.08 à une seconde, 0.21 à cinq secondes, 0.35 à quinze, 0.40 à trente et 0.48
à une minute, puis retombe à 0.40 à deux minutes et à 0.20 à cinq minutes. Cette baisse relève du bruit et non d'un renversement~: l'erreur type, calculée à partir de blocs de 10 minutes, passe de
0.07 tick à une minute à 0.15 tick à cinq minutes, soit un rapport de 2.28 contre $\sqrt5=2.24$. Le simulateur peut retirer le mouvement propre du prix efficient,
$\varepsilon_t(P^\ast_{t+h}-P^\ast_t)$, de l'impact de chaque transaction~; le reste se stabilise à 0.46 tick à partir d'une minute (0.48 à deux minutes, 0.46 à cinq
minutes, erreurs types de 0.04), soit 85~\% de la demi-fourchette effective, valeur proche de la vérité. La fourchette réalisée passe de 0.50 tick à 50 millisecondes à 0.056 à
une minute~: les fournisseurs de liquidité conservent un dixième de ce qu'ils facturent.

\begin{omfigure}
\begin{tikzpicture}
\begin{axis}[width=10cm,height=5.4cm,xmode=log,xlabel={\small horizon $h$ (secondes)},ylabel={\small ticks},tick label style={font=\footnotesize},
  xmin=0.035,xmax=420,ymin=-0.05,ymax=0.62,grid=major,grid style={gray!25},legend columns=2,legend style={font=\footnotesize,at={(0.5,-0.3)},
  anchor=north,draw=none,fill=none,/tikz/every even column/.append style={column sep=8pt}},table/col sep=comma]
\addplot[omDef,very thick,mark=*,error bars/.cd,y dir=both,y explicit] table[x=h,y=impact,y error=se]{figdata/microstructure/05-decomposing-the-spread/horizon.csv};
\addplot[omThm,very thick,mark=square*,error bars/.cd,y dir=both,y explicit] table[x=h,y=clean,y error=clean_se]{figdata/microstructure/05-decomposing-the-spread/horizon.csv};
\addplot[omProp,thick,mark=triangle*] table[x=h,y=realised]{figdata/microstructure/05-decomposing-the-spread/horizon.csv};
\addplot[black,dashed,domain=0.035:420] {0.4491};
\legend{impact,impact net du mouvement de $P^\ast$,demi-fourchette réalisée,part de sélection adverse réelle}
\end{axis}
\end{tikzpicture}

\omcaption{Impact sur le prix et demi-fourchette réalisée par horizon sur la journée simulée (pondéré par le nombre d'actions, avec des erreurs types issues de blocs de 10 minutes)~; la
demi-fourchette effective vaut 0.535~ticks. L'impact atteint la vraie part de sélection adverse vers une minute~; au-delà, les mouvements propres du prix efficient le rendent bruité. Données~:
\texttt{mx\_decomp.by\_horizon}.}\label{fig:mx:decomposing-the-spread:horizon}
\end{omfigure}

Ventiler les transactions selon l'indicateur du simulateur montre qui paie. À une minute, les transactions informées ont une demi-fourchette effective de 0.54~ticks, un impact de
2.11 et une demi-fourchette réalisée de $-1.57$~: chaque action informée coûte au fournisseur de liquidité un tick et demi. Les transactions des \omterm{def:mx:why-there-is-a-spread:informed}{noise traders} paient 0.53, déplacent le mid de 0.045 et laissent 0.49.
La faible marge du fournisseur de liquidité est la différence entre deux grands nombres.

\section{Décompositions structurelles}

Avant que les données de cotation intrajournalières ne soient courantes, les fourchettes étaient décomposées à partir des seuls prix de transaction, à l'aide d'un modèle de la manière dont les prix réagissent aux transactions.

Glosten et Harris (1988) écrivent le prix de transaction sous la forme $p_t=m_t+\varepsilon_t(c_0+c_1v_t)$, une partie transitoire linéaire en la taille de la transaction $v_t$, et
la révision du mid s'écrit $m_t-m_{t-1}=\varepsilon_t(z_0+z_1v_t)+y_t$, une partie permanente~; les variations de prix sont alors linéaires en $\Delta\varepsilon_t$,
$\Delta(\varepsilon_tv_t)$, $\varepsilon_t$ et $\varepsilon_tv_t$.

\begin{definition}[modèle de Huang--Stoll]\label{def:mx:decomposing-the-spread:hs}
Le \emph{modèle de Huang--Stoll}\index{modèle de Huang--Stoll} (forme de base) écrit la variation du prix de transaction entre deux transactions comme
\[
\Delta p_t=\frac S2\,\Delta\varepsilon_t+\lambda\,\frac S2\,\varepsilon_{t-1}+e_t,
\]
où $S$ est la fourchette des transactions et $\lambda$ la part de la demi-fourchette dont les cotations se déplacent après une transaction~: sélection adverse plus inventaire, ce que
la forme de base ne sait pas séparer. Ses extensions les séparent à l'aide d'un modèle de l'autocorrélation des signes des transactions.
\end{definition}

\begin{definition}[modèle de Madhavan--Richardson--Roomans]\label{def:mx:decomposing-the-spread:mrr}
Le \emph{modèle de Madhavan--Richardson--Roomans}\index{modèle de Madhavan--Richardson--Roomans} (MRR) laisse le prix efficient évoluer avec la surprise dans
le flux d'ordres et le prix de transaction osciller autour de lui~:
\[
\mu_t=\mu_{t-1}+\theta\left(x_t-\E[x_t\mid x_{t-1}]\right)+u_t,\qquad p_t=\mu_t+\phi\,x_t+\xi_t,
\]
avec $x_t$ le signe de la transaction, $\E[x_t\mid x_{t-1}]=\rho x_{t-1}$, $\theta$ le contenu en information d'une transaction et $\phi$ le coût qui ne dépend pas de
l'information. La fourchette implicite vaut $2(\phi+\theta)$ et la part d'information $\theta/(\phi+\theta)$.
\end{definition}

En différenciant, $\Delta p_t=(\phi+\theta)x_t-(\phi+\rho\theta)x_{t-1}+u_t+\Delta\xi_t$. Le MRR estime $(\theta,\phi,\rho)$ par la méthode des
moments généralisée (livre~4, chapitre~11)~; avec $\rho$ pris comme autocorrélation empirique des signes, les conditions de moments qui rendent l'erreur orthogonale
à $x_t$ et $x_{t-1}$ sont les équations normales d'une régression par moindres carrés, ce que fait la construction.

\omcode{code/firm/spreaddecomp/firm_spreaddecomp.py}{72}{81}{Le modèle MRR~: information et coût issus d'une régression sur les signes courant et précédent.}\label{lst:mx:decomposing-the-spread:mrr}

Le jour simulé, les trois modèles s'accordent entre eux et non avec la vérité. Huang--Stoll donne une fourchette des transactions de 1.07~ticks et
$\lambda=0.09$~; MRR une part d'information de 0.13 avec $\rho=0.33$~; Glosten--Harris une part de sélection adverse de 0.10 pour une transaction de 100~actions, qui baisse avec la taille
(les ordres informés du simulateur sont un peu plus petits, 109 actions en moyenne contre 120). Les trois modèles supposent que les cotations ont absorbé
l'information d'une transaction avant la transaction suivante. Sur le marché simulé, le mid se rapproche du prix efficient en quelques dizaines de secondes, par des mises à jour des cotations entre
les transactions, et un modèle construit sur les variations d'une transaction à l'autre ne voit que la première révision~: il prend le reste pour du bruit et sous-estime la sélection adverse
d'un facteur de six à neuf.

\section{L'approche par modèle vectoriel autorégressif}

Hasbrouck (1991) laisse les données choisir la dynamique. On prend la variation du mid-quote $r_t$ de juste avant la transaction $t$ à juste avant la transaction $t+1$ et le signe de la
transaction $x_t$, et on ajuste un modèle vectoriel autorégressif (VAR, livre~4, chapitre~20)~:
\[
r_t=\sum_{i=1}^{L}a_ir_{t-i}+\sum_{i=0}^{L}b_ix_{t-i}+v_{1,t},\qquad x_t=\sum_{i=1}^{L}c_ir_{t-i}+\sum_{i=1}^{L}d_ix_{t-i}+v_{2,t}.
\]
La transaction entre dans l'équation des cotations au retard zéro~: les cotations réagissent aux transactions à l'intérieur de l'intervalle, et non l'inverse. La réponse impulsionnelle cumulée du mid
à un choc unitaire de transaction, cumulée jusqu'à ce qu'elle cesse de varier, est l'impact permanent d'une transaction~: l'information qu'elle porte, quel que soit le chemin que suivent les cotations
pour l'absorber. Contrairement au MRR, il laisse l'absorption prendre de nombreuses transactions.

La méthode demande encore autant de transactions que le permettent les retards. Le jour simulé, l'impact permanent représente 0.14 de la demi-fourchette effective avec un retard, 0.28 avec
cinq, 0.38 avec dix, 0.49 avec vingt, 0.56 avec cinquante et 0.48 avec cent (\cref{fig:mx:decomposing-the-spread:methods})~: il croît tant que les retards
couvrent davantage l'ajustement, puis les coefficients supplémentaires ajoutent du bruit. Il n'atteint jamais la vérité, pour une seconde raison~: le VAR est linéaire en les signes,
alors que dans le simulateur les \omterm{def:mx:why-there-is-a-spread:informed}{traders informés} négocient lorsque l'écart $P^\ast-m$ est grand, une règle non linéaire que la réponse linéaire lisse en moyenne.

\begin{omfigure}
\begin{tikzpicture}
\begin{axis}[width=10cm,height=6.2cm,xbar,bar width=7pt,xlabel={\small part de sélection adverse de la demi-fourchette effective},
  tick label style={font=\footnotesize},label style={font=\small},xmin=0,xmax=1,ymin=-0.6,ymax=10.6,ytick={0,1,2,3,4,5,6,7,8,9,10},y dir=reverse,
  yticklabels={vérité,réalisée (1 min),Glosten--Harris (100),Huang--Stoll,MRR,VAR 1 retard,VAR 5 retards,VAR 10 retards,VAR 20 retards,VAR 50 retards,VAR 100 retards},
  grid=major,grid style={gray!25},table/col sep=comma]
\addplot[omDef,fill=omDef!40] table[x=share,y=k]{figdata/microstructure/05-decomposing-the-spread/estimators.csv};
\draw[black,thick,dashed] (axis cs:0.839,-0.6) -- (axis cs:0.839,10.6);
\end{axis}
\end{tikzpicture}

\omcaption{La part de sélection adverse de la fourchette selon chaque méthode le jour simulé, comparée à la vérité (en tirets), calculée à partir du prix efficient.
Seule la méthode fondée sur la demi-fourchette réalisée à une minute, qui voit les cotations, s'en approche. Données~: \texttt{mx\_decomp.estimators}, \texttt{var\_by\_lags}.}\label{fig:mx:decomposing-the-spread:methods}
\end{omfigure}

La leçon n'est pas que les modèles structurels aient tort~: chacun a raison sur le marché qu'il suppose. Elle tient à ce que la méthode doive correspondre à la manière dont les cotations absorbent
l'information sur le marché en question, et que, lorsqu'un simulateur à vérité connue existe, la méthode doit être évaluée sur lui avant d'utiliser son chiffre.

\section{Fourchettes à partir de données journalières}

\begin{definition}[estimateur de fourchette basse fréquence]\label{def:mx:decomposing-the-spread:lowfreq}
Un \emph{estimateur de fourchette basse fréquence}\index{estimateur de fourchette basse fréquence} déduit la fourchette effective des prix journaliers (ou par barres)~: clôtures, plus hauts et
plus bas, pour les marchés et les périodes sans cotations intrajournalières. L'estimateur de Roll (livre~4, chapitre~21) utilise l'autocovariance des variations de prix~;
Corwin et Schultz (2012) utilisent les amplitudes haut--bas~; Abdi et Ranaldo (2017), les clôtures et les milieux d'amplitude.
\end{definition}

Corwin et Schultz observent que le plus haut est en général un achat au cours vendeur et le plus bas une vente au cours acheteur, si bien qu'une amplitude journalière contient la fourchette une fois et la
volatilité une fois, et la partie volatilité croît avec la fenêtre, tandis que la partie fourchette ne croît pas. Comparer deux amplitudes d'une seule journée à l'amplitude sur les deux
jours les sépare. Avec $H$ et $L$ les logarithmes du plus haut et du plus bas,
\[
\beta=\sum_{j=0}^{1}(H_{t+j}-L_{t+j})^2,\quad \gamma=\bigl(\max(H_t,H_{t+1})-\min(L_t,L_{t+1})\bigr)^2,\quad
\alpha=\frac{\sqrt{2\beta}-\sqrt\beta}{3-2\sqrt2}-\sqrt{\frac{\gamma}{3-2\sqrt2}},
\]
et la fourchette relative vaut $S=2(e^\alpha-1)/(1+e^\alpha)$, mise à zéro lorsqu'elle est négative. Abdi et Ranaldo utilisent la clôture $c_t$ et le milieu d'amplitude
$\eta_t=\tfrac12(H_t+L_t)$~: $S^2=4\,\E[(c_t-\eta_t)(c_t-\eta_{t+1})]$.

Les deux méthodes exigent que la fourchette représente une part visible de l'amplitude. Le jour simulé, la fourchette effective vaut 1.07~points de base du prix. Sur des barres de 30
secondes, une minute, cinq minutes et quinze minutes (761, 390, 78 et 26 barres), Corwin--Schultz donne 0.5, 0.6, 2.0 et 2.6 points de base et Abdi--Ranaldo
0.1, 0.0, 3.3 et 0.0~: nulle part fiables, car sur un marché liquide avec une fourchette d'un tick, l'essentiel de l'amplitude relève de la volatilité, et les estimateurs soustraient
deux grands nombres. Ils ont été conçus pour des données journalières sur des actions moins liquides, dont les fourchettes représentent une plus grande part de l'amplitude, et il est préférable de les vérifier, comme ici,
sur une période où la fourchette effective est connue.

\section{Tutoriel~: qui a payé les deux centimes}\label{tut:mx:decomposing-the-spread:tut}

\begin{tutorial}
\textbf{Objectif.} Décomposer la fourchette du jour simulé par chaque méthode et évaluer chacune par rapport à la vérité.
\textbf{État final~:} les \cref{fig:mx:decomposing-the-spread:horizon,fig:mx:decomposing-the-spread:methods} et les chiffres des sections 2 à 5.
\begin{tutsteps}
\item \textbf{La journée.} \texttt{mx\_decomp.day()}~: la séance de 6.5 heures de \texttt{firm.tape}, avec son prix efficient \texttt{v} et
l'indicateur \texttt{informed} de chaque transaction.
\item \textbf{Vérité.} \texttt{truth(tape)}~: la demi-fourchette effective et $\E[\varepsilon(P^\ast-m)]$, globalement et par type de trader.
\item \textbf{Horizons.} \texttt{by\_horizon(tape)}~: \texttt{decompose} à huit horizons, avec les erreurs types par blocs et l'impact net du mouvement de $P^\ast$.
\item \textbf{Modèles.} \texttt{estimators(tape)}~: Glosten--Harris, Huang--Stoll, MRR et le VAR~; \texttt{var\_by\_lags(tape)} pour le balayage des retards.
\item \textbf{Données journalières.} \texttt{low\_frequency(tape, bar\_s)} pour des barres de 30 secondes à 15 minutes~; tracer avec \texttt{fig\_decomp.py}.
\end{tutsteps}
\textbf{Pour aller plus loin.} Rendre les fournisseurs de liquidité de \texttt{firm.tape} plus lents à suivre le prix efficient (un taux de mise à jour des cotations plus faible)
et observer l'horizon auquel l'impact se stabilise reculer~; lancer la décomposition sur la journée enregistrée de \texttt{firm.exchsim}, où les agents de tape du livre~7 négocient
contre le \omterm{def:mx:the-limit-order-book:engine}{moteur d'appariement}.
\end{tutorial}

\section{Construction~: la décomposition de la fourchette}\label{bld:mx:decomposing-the-spread:spreaddecomp}

\begin{build}
\bfield{Rôle} Les mesures et \omterm{def:mx:decomposing-the-spread:decomposition}{décompositions de la fourchette} utilisées par l'analyse des coûts de transaction du chapitre~19, par les indicateurs de qualité de marché du chapitre~24
et par les études de mark-out du livre~11.
\bfield{Interface} \texttt{mid\_at}, \texttt{quoted\_spread}, \texttt{decompose(trades, times, mids, h)}~; \texttt{glosten\_harris}, \texttt{huang\_stoll},
\texttt{mrr}~; \texttt{hasbrouck\_var(r, x, lags, steps)}~; \texttt{corwin\_schultz}, \texttt{abdi\_ranaldo} (l'estimateur de Roll se trouve dans
\texttt{firm.spreadmodels}).
\bfield{Règles} Les transactions portent le signe de l'agresseur~; le mid est celui en vigueur juste avant la transaction~; les moyennes sont pondérées par les actions~; les modèles structurels
sont ajustés par moindres carrés avec numpy~; les estimations négatives de fourchette basse fréquence sont ramenées à zéro, comme le font les articles d'origine.
\bfield{Tests d'acceptation} \texttt{code/firm/spreaddecomp/tests/}~: l'effective égale la réalisée plus l'impact sur les transactions saisies à la main~; MRR, Huang--Stoll et
Glosten--Harris retrouvent les paramètres de marchés simulés à partir de leurs propres modèles~; le VAR retrouve un impact permanent introduit volontairement~; Corwin--Schultz et
Abdi--Ranaldo retrouvent une fourchette de 1~\% introduite volontairement à partir des cours à la minute.
\bfield{Pour aller plus loin} La décomposition de Huang--Stoll à trois termes avec autocorrélation du signe~; MRR par MMG complète avec erreurs types~; la part d'information de Hasbrouck
entre plateformes.
\end{build}

\begin{omsources}
\item L.~R.~Glosten et L.~E.~Harris, ``Estimating the components of the bid/ask spread'', \emph{Journal of Financial Economics} 21(1), 1988.
\item R.~D.~Huang et H.~R.~Stoll, ``The components of the bid-ask spread: a general approach'', \emph{Review of Financial Studies} 10(4), 1997.
\item A.~Madhavan, M.~Richardson et M.~Roomans, ``Why do security prices change? A transaction-level analysis of NYSE stocks'', \emph{Review of Financial
Studies} 10(4), 1997.
\item J.~Hasbrouck, ``Measuring the information content of stock trades'', \emph{Journal of Finance} 46(1), 1991.
\item S.~A.~Corwin et P.~Schultz, ``A simple way to estimate bid-ask spreads from daily high and low prices'', \emph{Journal of Finance} 67(2), 2012.
\item F.~Abdi et A.~Ranaldo, ``A simple estimation of bid-ask spreads from daily close, high, and low prices'', \emph{Review of Financial Studies} 30(12), 2017.
\item U.S. Securities and Exchange Commission, \emph{Disclosure of Order Execution Information}, Release 34-99679, 89 FR 26428, 2024~; ordonnance d'exemption,
Release 34-105136, 2026.
\end{omsources}

\section{Exercices}

\begin{exercise}[$\star$]\label{exo:mx:decomposing-the-spread:1}
Le marché affiche 100.00 au cours acheteur et 100.04 au cours vendeur. Un achat de 1\,000 actions s'exécute à 100.03, et une minute plus tard le mid vaut 100.025. Calculer les demi-fourchettes affichée et effective,
l'amélioration de prix, la demi-fourchette réalisée et l'impact par action, ainsi que la fourchette effective en points de base.
\end{exercise}

\begin{exercise}[$\star$]\label{exo:mx:decomposing-the-spread:2}
Une vente s'exécute à 49.98 avec un mid à 50.00~; cinq minutes plus tard, le mid vaut 49.97. Calculer les demi-fourchettes effective et réalisée ainsi que l'impact, et expliquer
ce que signifie une fourchette réalisée négative pour le fournisseur de liquidité.
\end{exercise}

\begin{exercise}[$\star$]\label{exo:mx:decomposing-the-spread:3}
Le \omterm{def:mx:decomposing-the-spread:hs}{modèle de Huang--Stoll} de base renvoie une fourchette négociée de 0.04 et $\lambda=0.6$. De combien les cotations bougent-elles après un achat, et quelle part de la demi-fourchette
reste-t-il pour couvrir les coûts de traitement~?
\end{exercise}

\begin{exercise}[$\star\star$]\label{exo:mx:decomposing-the-spread:4}
Un ajustement MRR donne $\theta=0.012$, $\phi=0.008$ et $\rho=0.3$. Calculer la fourchette implicite et la part d'information, ainsi que la variation de prix attendue pour un achat qui
suit un achat et pour un achat qui suit une vente. Pourquoi diffèrent-elles~?
\end{exercise}

\begin{exercise}[$\star\star$]\label{exo:mx:decomposing-the-spread:5}
Sur la journée simulée, l'impact vaut 0.48 tick à une minute et 0.20 à cinq minutes, avec des erreurs types de 0.07 et 0.15. Le mouvement à une minute a-t-il été inversé~? Que nous apprend
le rapport des erreurs types, et que montre l'impact net du mouvement de $P^\ast$~?
\end{exercise}

\begin{exercise}[$\star\star$]\label{exo:mx:decomposing-the-spread:6}
Deux jours consécutifs ont des plus hauts de 100.10 et 100.12 et des plus bas de 99.94 et 99.96. Calculer $\beta$, $\gamma$ et la fourchette de Corwin--Schultz.
\end{exercise}

\begin{exercise}[$\star\star\star$]\label{exo:mx:decomposing-the-spread:7}
\emph{Programmation.} Sur la journée simulée, séparer les transactions selon l'indicateur \texttt{informed} et décomposer chaque groupe à une minute. Qui paie la marge du fournisseur,
et comment la moyenne pondérée par les actions retrouve-t-elle la fourchette réalisée globale de 0.056 tick~?
\end{exercise}

\begin{exercise}[$\star\star\star$]\label{exo:mx:decomposing-the-spread:8}
\emph{Trouver la faille.} «~Les ordres enfants de notre algorithme affichent une fourchette réalisée négative à cinq minutes pour les teneurs de marché qui les exécutent~; notre flux est donc
informé, et nous méritons de meilleurs prix.~»
\end{exercise}

\section{Problème~: qui a payé les deux cents~?}

\begin{problem}[{Problème du week-end --- qui a payé les deux cents~?}]\label{pb:mx:decomposing-the-spread:1}
Un fournisseur de liquidité veut savoir quelle part de sa fourchette paie le fait de se faire prendre au piège, et un régulateur veut savoir à quel horizon publier. Sur le
marché simulé, la vérité est connue~; évaluer chaque méthode par rapport à elle.

\medskip\noindent\textbf{Partie I --- Fourchettes et impact.}
\begin{enumerate}
\item Écrire l'identité entre demi-fourchettes effective, réalisée et d'impact, et la démontrer.
\item Quelles sont la \omterm{def:mx:decomposing-the-spread:quoted}{fourchette affichée} et la demi-fourchette effective pondérée par les actions de la journée simulée~?
\item Définir la \omterm{def:mx:decomposing-the-spread:adverse}{composante de sélection adverse} et calculer la vérité~: en ticks et en part.
\item À quelle distance le prix efficient se situe-t-il du mid lorsque les \omterm{def:mx:why-there-is-a-spread:informed}{traders informés} et les \omterm{def:mx:why-there-is-a-spread:informed}{noise traders} négocient~?
\item Donner l'impact à 50 millisecondes, à une seconde, à quinze secondes et à une minute.
\end{enumerate}

\medskip\noindent\textbf{Partie II --- L'horizon.}
\begin{enumerate}[resume]
\item Énoncer et démontrer la proposition de convergence.
\item Pourquoi l'impact baisse-t-il entre une et cinq minutes~?
\item Quel est l'impact net du mouvement du prix efficient, et pourquoi seule une simulation peut-elle le calculer~?
\item Donner la fourchette réalisée à 50 millisecondes et à une minute, et dire ce que garde le fournisseur.
\item Lequel des cinq horizons de la Rule 605 est le plus proche de l'horizon où l'impact se stabilise ici~?
\end{enumerate}

\medskip\noindent\textbf{Partie III --- Modèles.}
\begin{enumerate}[resume]
\item Dériver l'équation MRR de la variation de prix.
\item Que renvoient Huang--Stoll, MRR et Glosten--Harris~?
\item Pourquoi sous-estiment-ils la sélection adverse sur ce marché~?
\item Que renvoie le VAR avec 1, 5, 20, 50 et 100 retards, et pourquoi la courbe n'est-elle pas monotone~?
\item Pourquoi un VAR linéaire ne peut-il pas atteindre la vérité ici~?
\end{enumerate}

\medskip\noindent\textbf{Partie IV --- Données journalières et verdict.}
\begin{enumerate}[resume]
\item Expliquer l'idée de l'estimateur de Corwin--Schultz.
\item Que renvoient Corwin--Schultz et Abdi--Ranaldo sur des barres d'une minute et de cinq minutes, par rapport à la fourchette effective~?
\item Quand peut-on faire confiance aux estimateurs basse fréquence~?
\item Énoncer le \emph{résultat nommé}~: la part de sélection adverse estimée par chaque méthode par rapport à la vérité du simulateur, et l'horizon auquel la
fourchette réalisée se stabilise.
\item En une phrase~: qui a payé les deux cents~?
\end{enumerate}
\end{problem}

\section{Questions d'entretien}

\begin{interviewq}[$\star$ \iqroles{trader, chercheur}]\label{iq:mx:decomposing-the-spread:1}
Définissez les fourchettes effective et réalisée. Que mesure la différence entre les deux~?
\end{interviewq}

\begin{interviewq}[$\star\star$ \iqroles{chercheur}]\label{iq:mx:decomposing-the-spread:2}
Comment choisissez-vous l'horizon d'une fourchette réalisée, et que se passe-t-il s'il est trop court ou trop long~?
\end{interviewq}

\begin{interviewq}[$\star\star$ \iqroles{chercheur}]\label{iq:mx:decomposing-the-spread:3}
Écrivez le modèle MRR et expliquez comment ses deux paramètres sont identifiés.
\end{interviewq}

\begin{interviewq}[$\star\star$ \iqroles{trader}]\label{iq:mx:decomposing-the-spread:4}
Votre desk de teneur de marché dégage une fourchette réalisée positive à une seconde mais perd de l'argent. Comment est-ce possible~?
\end{interviewq}

\begin{interviewq}[$\star\star$ \iqroles{chercheur}]\label{iq:mx:decomposing-the-spread:5}
Que mesure le VAR de Hasbrouck, et pourquoi la transaction entre-t-elle dans l'équation de cotation au retard zéro~?
\end{interviewq}

\begin{interviewq}[$\star\star\star$ \iqroles{chercheur}]\label{iq:mx:decomposing-the-spread:6}
Vous avez besoin de fourchettes pour un échantillon d'actions des années 1970 avec seulement des données journalières. Qu'utilisez-vous, et comment le vérifiez-vous~?
\end{interviewq}
